\begin{frame}{예 2: $\gamma$ 를 올리면 답이 뒤집힌다}
  $\gamma = 0.9$ 로 두 정책의 가치를 반복 대입(3.3절 벨만방정식
  고정점)으로 풀어보면:
  \[
  V^A = [-63.36,\; -71.51,\; -100], \qquad
  V^B(0) = \frac{1}{1-\gamma^2} = \frac{1}{0.19} \approx 5.26
  \]
  \begin{itemize}
    \item 벨만방정식 세 등호: $V(0) = 1 + 0.9\,V(1)$,
      $V(1) = 2 + 0.9\,(0.5\,V(2) + 0.5\,V(0))$, $V(2) = -10$
      (코드로 검산: $V^A(0) = -63.3613$).
    \item $V^A(2) = -10/(1-0.9) = -100$ --- ``한 번 $-10$ 루프에
      빠지면 그 상태의 가치는 $-100$''.
    \item $\gamma = 0.9$ 에서는 \textbf{정책 A(−63.36)가
      정책 B(5.26)보다 훨씬 나쁘다} --- ``50\% 확률로 $-100$ 짜리
      구덩이''가 2스텝짜리 $+2$ 보상의 기대를 훨씬 넘는다.
  \end{itemize}
  \vskip0.5em
  그런데 \textbf{$\gamma$ 를 낮추면 답이 뒤집힌다}: $\gamma = 0.1$
  이면 $V^A(0) = 1.15 > V^B(0) = 1.01$ --- \kb{정책 A가 오히려
  좋다}. $\gamma \to 0$ 극한에서는 ``앞으로 1스텝만'' 보므로 상태 1에서
  $+2$ 를 즉시 주는 A가 불리하지 않다.
  \vskip0.3em
  {\footnotesize 그림: 오른쪽 패널에서 두 정책의 상태 0 가치 곡선이
  크로스오버 $\gamma^* \approx 0.27$ 에서 교차.}
\end{frame}
